\documentclass[10pt,a4paper]{article}
\usepackage[margin=0.58in]{geometry}
\usepackage[T1]{fontenc}
\usepackage{lmodern}
\usepackage{microtype}
\usepackage{enumitem}
\usepackage[hidelinks]{hyperref}
\usepackage{xcolor}
\usepackage{titlesec}
\usepackage{tabularx}
\usepackage{array}
\usepackage{parskip}

\definecolor{accent}{HTML}{1F1F1F}
\pagestyle{empty}
\setlength{\parindent}{0pt}
\setlength{\parskip}{0pt}
\setlist[itemize]{leftmargin=1.18em,itemsep=1.1pt,topsep=1.5pt,parsep=0pt,partopsep=0pt}
\titleformat{\section}{\large\bfseries\color{accent}}{}{0em}{}[\vspace{-0.40em}\rule{\textwidth}{0.32pt}]
\titlespacing*{\section}{0pt}{0.55em}{0.28em}

\newcommand{\entry}[4]{%
  \begin{tabularx}{\textwidth}{@{}X>{\raggedleft\arraybackslash}p{0.28\textwidth}@{}}
    \textbf{#1} & #2 \\
    \textit{#3} & \textit{#4}
  \end{tabularx}
}
\newcommand{\project}[4]{%
  \begin{tabularx}{\textwidth}{@{}X>{\raggedleft\arraybackslash}p{0.23\textwidth}@{}}
    \textbf{#1} --- \textit{#2} & #3 \\
    \multicolumn{2}{@{}l@{}}{\footnotesize #4}
  \end{tabularx}
}

\begin{document}

{\LARGE\bfseries Zhiheng Zhang}\\[-0.05em]
\href{mailto:m15601654187@163.com}{m15601654187@163.com}
\quad $\vert$ \quad
\href{mailto:zhihezhang@student.unimelb.edu.au}{zhihezhang@student.unimelb.edu.au}\\
\href{https://github.com/zhiheng-zhang-Mera}{github.com/zhiheng-zhang-Mera}
\hfill Melbourne, Australia

\vspace{0.30em}

\section{Research Profile}
Computer Science researcher and research engineer focused on \textbf{trustworthy AI systems, scientific agents, long-horizon orchestration, empirical evaluation, and reproducible AI-for-science workflows}.

\section{Education}
\entry{University of Melbourne}{Melbourne, Australia}{Master's degree in Computer Science}{Expected 2026}
\entry{University of British Columbia}{Vancouver, Canada}{Bachelor of Science in Computer Science}{Completed Summer 2024}

\section{Selected Research and Engineering Projects}

\project{DS-Hns}{Long-horizon autonomous software-work runtime}{2026--present}{\href{https://github.com/zhiheng-zhang-Mera/DS-Hns}{github.com/zhiheng-zhang-Mera/DS-Hns}}
\begin{itemize}
  \item Building an execution layer for unattended software work around the official DeepSeek Harness interface, with optional capabilities isolated as extensions.
  \item Implemented ordered/scheduled task queues, hardware-adaptive local concurrency, unified task lifecycle and terminal events, task history, and terminal notifications.
  \item Designed extension-failure isolation, restart/recovery-oriented execution, and an idempotent Windows installer with explicit optional-plugin flow.
  \item Using the system as a testbed for \textbf{reliable coding agents, persistent execution, recovery, scheduling, and completion verification}.
\end{itemize}

\project{Codex Boss}{Multi-agent orchestration and evidence-based acceptance platform}{2026--present}{\href{https://github.com/zhiheng-zhang-Mera/Codex-Boss}{github.com/zhiheng-zhang-Mera/Codex-Boss}}
\begin{itemize}
  \item Building a platform for long-running AI workflows in which multiple agents can plan, execute, critique, and adjudicate work while preserving durable task state.
  \item Developed platform foundations spanning durable state/events, capability boundaries, plugin isolation, knowledge/data lifecycle, dogfooding, and acceptance/verification.
  \item Investigating how autonomous systems can distinguish genuine task completion from superficially plausible but weak evidence.
  \item Current semantic-acceptance work deliberately retains unresolved failure cases: a vacuous case is rejected, while a meaningful case still exposes a false negative in the evidence reader rather than being reported as solved.
  \item Current regression evidence on the relevant branch includes \textbf{2,508 unit tests across 215 files}, plus Electron, renderer, and test typechecking.
\end{itemize}

\project{Quant-Ultra}{Empirical quantitative-research and decision-system project}{2026--present}{\href{https://github.com/zhiheng-zhang-Mera/Quant-ultra}{github.com/zhiheng-zhang-Mera/Quant-ultra}}
\begin{itemize}
  \item Developed a personal quantitative-research platform with emphasis on robustness, temporal evaluation, data integrity, risk-aware decision making, and reproducible workflows.
  \item Shifted the project away from headline-return claims toward evaluation quality, non-stationary data, temporal leakage, walk-forward reasoning, and auditable experiment infrastructure.
\end{itemize}

\project{Privacy Lens}{Trustworthy-software / privacy project}{}{}
\begin{itemize}
  \item Explored privacy-risk interpretation and user-facing warning mechanisms with attention to bounded interpretation, auditability, evidence boundaries, and trustworthy software behavior.
  \item Supports broader interests in trustworthy AI, software governance, and high-stakes automated systems.
\end{itemize}


\section{Research Interests}
Foundation models and agents for science; trustworthy and auditable AI; long-horizon agent orchestration; simulation/data-analysis workflows; uncertainty-aware evaluation; reproducible AI-for-science systems.

\section{Technical and Research Strengths}
\begin{itemize}
  \item End-to-end research engineering: problem formulation, system design, iterative implementation, testing, debugging, and empirical evaluation.
  \item Multi-module software design with explicit persistence, lifecycle, extensibility, isolation, recovery, and acceptance criteria.
  \item Git/GitHub-based experimental development, branch-based iteration, regression testing, and reproducible project documentation.
  \item Practical multi-AI workflow design: using multiple AI systems as coding/research collaborators while independently reviewing plans, evidence, failure modes, and completion claims.
\end{itemize}

\section{Current Research Direction}
I am interested in \textbf{trustworthy scientific agents and reproducible AI-for-science workflows}: how agents should combine tools, simulation/data-analysis steps, persistent state, uncertainty/evidence checks, and explicit acceptance criteria so that long-running scientific workflows remain auditable rather than merely plausible.

\end{document}
